\section{$D_3$：从六个群元到三角形的对称模式}
\label{sec:d3-comprehensive}

正三角形的旋转与反射给出六阶群。固定这两个生成元后，可以依次算出乘法、子群、群作用和表示，再用特征标求投影与张量积分解。
以下例题沿用同一组群元和表示矩阵；使用一般定理时，同时检查其假设在这个六阶群中是否满足。
约定置换从右向左复合，$D_3$ 表示\emph{六阶}二面体群。

\begin{workedexample}[由生成元关系计算乘法]
\label{ex:d3-comprehensive-multiplication}
取
\[
G=D_3=\langle r,s\mid r^3=e,\ s^2=e,\ srs=r^{-1}\rangle.
\]
求群元的正规形、乘法和元素阶，并验证这套运算满足群公理。

\textbf{解。} 关系 $sr=r^{-1}s$ 把任意词中的 $s$ 移到右端，再约去 $r^3$ 与 $s^2$，
所以每个群元均可写成 $r^a s^b$，其中 $a\in\mathbb Z/3\mathbb Z$、$b\in\mathbb Z/2\mathbb Z$。
令 $r=(123)$、$s=(23)$，六个这样的词在 $S_3$ 中互不相同。
因此这些正规形恰好给出六个群元，展示所定义的群与 $S_3$ 同构。

把正规形写成 $(a,b)$，乘法成为
\[
(a,b)(c,d)=\bigl(a+(-1)^b c,\ b+d\bigr),
\qquad
(a,b)^{-1}=\bigl(-(-1)^b a,\ b\bigr).
\]
第一个坐标模 $3$、第二个坐标模 $2$。
为检查结合律，取第三个元素 $(f,h)$。左右两种三重乘法的第一个坐标同为
\[
a+(-1)^b c+(-1)^{b+d}f,
\]
第二个坐标同为 $b+d+h$，从而直接检查了结合律。
$(0,0)$ 是单位元，给出的逆元与 $(a,b)$ 左右相乘均为 $(0,0)$。
因此这套六元素运算满足\Cref{def:group}的全部条件。

以下表格每一格均由该公式计算，行元在左、列元在右：
\[
\begin{array}{c|rrrrrr}
 &e&r&r^2&s&rs&r^2s\\\hline
e&e&r&r^2&s&rs&r^2s\\
r&r&r^2&e&rs&r^2s&s\\
r^2&r^2&e&r&r^2s&s&rs\\
s&s&r^2s&rs&e&r^2&r\\
rs&rs&s&r^2s&r&e&r^2\\
r^2s&r^2s&rs&s&r^2&r&e
\end{array}
\]
特别地，$sr=r^2s\ne rs$，故 $G$ 非 Abel。
元素阶为
\[
\operatorname{ord}(e)=1,\quad
\operatorname{ord}(r)=\operatorname{ord}(r^2)=3,\quad
\operatorname{ord}(r^a s)=2.
\]
没有六阶元，所以六阶的 $G$ 并非循环群；$r,s$ 则确实共同生成整个 $G$。
\end{workedexample}

\begin{workedexample}[全部子群与子群格中的运算]
\label{ex:d3-comprehensive-subgroups}
列出全部子群，计算不同子群的交与共同生成的子群。

\textbf{解。} 由\Cref{thm:lagrange}，子群阶只能是 $1,2,3,6$。
阶为 $2$ 或 $3$ 的群为素数阶循环群；结合刚算出的元素阶，全部子群就是
\[
\{e\},\qquad R=\langle r\rangle=\{e,r,r^2\},\qquad
H_a=\langle r^a s\rangle=\{e,r^a s\}\ (a=0,1,2),\qquad G.
\]
阶为 $1$ 与 $6$ 的子群已确定，阶为 $2$ 与 $3$ 的子群又必须由相应阶的元素生成，所以此表没有遗漏。
对 $a\ne b$，不同反射相乘给出非平凡旋转：
\[
(r^a s)(r^b s)=r^{a-b}.
\]
所以
\[
H_a\cap H_b=\{e\},\quad \langle H_a,H_b\rangle=G,
\qquad R\cap H_a=\{e\},\quad \langle R,H_a\rangle=G.
\]
另一方面，集合 $H_0H_1=\{e,s,rs,r^2\}$ 只有四个元素，不是子群。
例如它包含 $r^2$ 却不包含 $(r^2)^2=r$。
因此“子群的交仍为子群”不能误改成“子群的集合乘积总是子群”。
这些交与生成计算已经确定了整个子群格：除顶端 $G$ 和底端 $\{e\}$ 外，
$R,H_0,H_1,H_2$ 之间没有包含关系，且均为极大子群。
\end{workedexample}

\begin{workedexample}[左右陪集、指标与不能定义的商群]
\label{ex:d3-comprehensive-cosets}
计算 $H_0$ 与 $R$ 的左右陪集，并检查陪集乘法是否依赖代表元。

\textbf{解。} 对 $H_0=\{e,s\}$，左陪集与右陪集分别为
\begin{align*}
H_0&=\{e,s\},&rH_0&=\{r,rs\},&r^2H_0&=\{r^2,r^2s\},\\
H_0&=\{e,s\},&H_0r&=\{r,r^2s\},&H_0r^2&=\{r^2,rs\}.
\end{align*}
每一行都把六个元素分成三个互不相交的二元素集合，因而
\[
[G:H_0]=3,\qquad |G|=|H_0|[G:H_0]=2\cdot3.
\]
左右陪集数相等，但 $rH_0\ne H_0r$；相等的是\emph{数目}，不是逐个集合。
更具体地，$H_0=eH_0=sH_0$，若试图定义左陪集乘法，则
\[
(eH_0)(rH_0)=rH_0,\qquad
(sH_0)(rH_0)=srH_0=r^2H_0.
\]
同一陪集更换代表元后得到不同输出，故 $G/H_0$ 只是陪集集合，不能按此定义商群。

对 $R$ 则有
\[
G=R\sqcup sR=R\sqcup Rs,\qquad
sR=Rs=\{s,rs,r^2s\},\qquad [G:R]=2.
\]
这里左右陪集相同，下一例题会从共轭直接确认 $R$ 的正规性。
\end{workedexample}

\begin{workedexample}[共轭类、中心化子与正规化子逐项计算]
\label{ex:d3-comprehensive-conjugacy}
由生成元的共轭作用求共轭类、中心化子和正规化子。

\textbf{解。} 由展示关系，$srs^{-1}=r^2$，而
\[
r(r^a s)r^{-1}=r^{a+2}s,\qquad
s(r^a s)s^{-1}=r^{-a}s.
\]
因此全部共轭类及类方程为
\[
\{e\},\qquad \{r,r^2\},\qquad \{s,rs,r^2s\},
\qquad 6=1+2+3.
\]
非平凡旋转与全部旋转交换，却不与反射交换，故
\[
C_G(r)=C_G(r^2)=R.
\]
两个不同反射不交换：其两种乘积分别是 $r^{a-b}$ 和 $r^{b-a}$，
当 $a\ne b$ 时二者不同。因此
\[
C_G(r^a s)=H_a,\qquad C_G(e)=G,
\qquad Z(G)=R\cap H_0=\{e\}.
\]
\Cref{thm:class-size}现在逐项给出 $6/3=2$ 与 $6/2=3$，恰与列出的类大小一致。

正规化子要求保持\emph{子群}，不要求固定其中每个元素。
共轭把 $r$ 变为 $r$ 或 $r^2$，所以 $N_G(R)=G$，而 $C_G(R)=R$。
对 $H_a$，保持该二阶子群就必须固定其唯一非单位元，故
\[
N_G(H_a)=C_G(H_a)=H_a.
\]
再加 $N_G(\{e\})=N_G(G)=G$、$C_G(\{e\})=G$、$C_G(G)=Z(G)=\{e\}$，
全部子群的正规化子与中心化子均已确定。
特别地，正规子群只有 $\{e\},R,G$；三个反射子群虽然两两共轭同构，却都不正规。
\end{workedexample}

\begin{workedexample}[从核与像读出第一、第二同构定理]
\label{ex:d3-comprehensive-isomorphism}
构造一个核为旋转子群 $R$ 的满同态，并在 $H_0$ 与 $R$ 上写出第一、第二同构定理的具体对应。

\textbf{解。} 用是否含反射区分两类群元，定义
\[
\varepsilon:G\longrightarrow\{1,-1\},\qquad
\varepsilon(r^a s^b)=(-1)^b.
\]
乘法公式给出 $\varepsilon(gh)=\varepsilon(g)\varepsilon(h)$，且
\[
\ker\varepsilon=R,\qquad\operatorname{im}\varepsilon=C_2.
\]
所以\Cref{thm:first-isomorphism}在这里把两类陪集具体识别为
\[
G/R=\{R,sR\}\cong C_2,\qquad R\mapsto1,\quad sR\mapsto-1.
\]
顶点置换同态 $G\to S_3$ 则核为 $\{e\}$、像为整个 $S_3$，重新给出 $G\cong S_3$。
也可以把旋转子群看成 $\mathbb Z\to R$ 的像：$n\mapsto r^n$ 的核是 $3\mathbb Z$，
故 $R\cong\mathbb Z/3\mathbb Z$。

对第二同构定理，取 $H=H_0$、正规子群 $N=R$。
因为 $H_0R=G$、$H_0\cap R=\{e\}$，自然映射为
\[
H_0\longrightarrow G/R,\qquad e\mapsto R,\quad s\mapsto sR.
\]
它是双射同态，因此
\[
H_0/(H_0\cap R)=H_0\cong H_0R/R=G/R.
\]
同一条二阶群结构在 $G$ 内由反射给出，在商群内由旋转陪集给出；
两者同构不表示它们是同一个集合。
\end{workedexample}

\begin{workedexample}[第三同构定理、对应定理与假设边界]
\label{ex:d3-comprehensive-correspondence}
将第三同构定理和对应定理代入 $G$ 的正规子群链，说明本例能得到哪些商群和子群对应。

\textbf{解。} 第三同构定理需要两个嵌套的正规子群。
本例全部正规子群已经穷举为 $\{e\},R,G$，因此唯一严格三级链为
\[
\{e\}\lhd R\lhd G.
\]
取 $N=\{e\}$、$K=R$，第三同构定理给出
\[
(G/N)/(K/N)\cong G/K=G/R\cong C_2.
\]
这是合法实例，却没有两个非平凡真正规子群 $N\subsetneq K$。
不能凭空把某个 $H_a$ 填入该链：它不正规，所需商群本就没有定义。

对应定理应用于投影 $\pi:G\to G/R$ 时尤其容易核查。
$G/R=C_2$ 只有 $\{R\}$ 和 $G/R$ 两个子群，其逆像分别为
\[
\pi^{-1}(\{R\})=R,\qquad \pi^{-1}(G/R)=G.
\]
这恰好列出 $G$ 中所有\emph{包含 $R$}的子群。
$H_a$ 不包含 $R$，所以它们不属于这次对应，而不是被定理遗漏。
\end{workedexample}

\begin{workedexample}[自同构、内自同构、特征子群与 $N/C$]
\label{ex:d3-comprehensive-automorphisms}
求全部自同构、内自同构和特征子群，再计算 $N/C$ 定理中的共轭映射。

\textbf{解。} 自同构保持元素阶，所以 $r$ 只能映到两个三阶元之一，$s$ 只能映到三个二阶元之一。于是必有
\[
\alpha_{u,v}(r)=r^u,\qquad\alpha_{u,v}(s)=r^v s,
\qquad u\in\{1,2\},\quad v\in\mathbb Z/3\mathbb Z.
\]
这些像都满足原展示关系，并仍生成 $G$，故全部六个选择均为自同构。
复合直接算成
\[
\alpha_{u,v}\alpha_{w,z}=\alpha_{uw,\,v+uz}.
\]
共轭 $c_r$ 对应 $(u,v)=(1,2)$，共轭 $c_s$ 对应 $(2,0)$；
它们生成这六个自同构。因此
\[
\operatorname{Aut}(G)=\operatorname{Inn}(G)\cong G/Z(G)\cong S_3,
\qquad\operatorname{Out}(G)=\{e\}.
\]
所有自同构保持唯一三阶子群 $R$，并置换三个 $H_a$。
特征子群必须先为正规子群，所以全部特征子群也恰为 $\{e\},R,G$。

把 $N/C$ 定理应用到 $R$，共轭同态为
\[
G=N_G(R)\longrightarrow\operatorname{Aut}(R),\qquad
r\mapsto\operatorname{id},\quad s\mapsto(r\mapsto r^{-1}).
\]
核是 $C_G(R)=R$，像是 $\operatorname{Aut}(C_3)=C_2$，所以
$N_G(R)/C_G(R)\cong C_2$。
应用到 $H_a$ 则 $N_G(H_a)/C_G(H_a)=\{e\}$，与 $\operatorname{Aut}(C_2)=\{e\}$ 相符。
可见正规化子商给出的是\emph{由环境群共轭实现}的自同构，不能一般地直接写成全部自同构。
\end{workedexample}

\begin{workedexample}[Sylow、Frattini 与两种不同的交]
\label{ex:d3-comprehensive-sylow}
确定全部 Sylow 子群，并检验不动点定理和 Frattini 论证的适用条件；另求所有极大子群的交。

\textbf{解。} 素因子 $2,3$ 的阶元分别已经由 $s,r$ 给出，满足 Cauchy 定理的存在结论。
Sylow 定理给出
\[
n_3\mid2,\quad n_3\equiv1\pmod3\ \Longrightarrow n_3=1,
\]
唯一 Sylow $3$-子群即 $R$。
对素数 $2$，$n_2\mid3$、$n_2\equiv1\pmod2$ 允许 $1$ 或 $3$；
定理的数论约束\emph{本身}尚不能在两者间选择。
已经列出的 $H_0,H_1,H_2$ 确认 $n_2=3$，共轭公式则直接验证 Sylow 共轭性。

$p$-群不动点定理也能在顶点集合 $X=\{1,2,3\}$ 上检查，仍取 $r=(123)$、$s=(23)$。
限制到 $R$ 时，唯一轨道大小是 $3$，因此 $|X|\equiv|X^R|=0\pmod3$；
限制到 $H_0$ 时，轨道是 $\{1\}$ 与 $\{2,3\}$，因此
$|X|=3\equiv |X^{H_0}|=1\pmod2$。
但 $G$ 本身不是 $p$-群，故不能用该定理断言 $Z(G)$ 非平凡；
前面已算出的 $Z(G)=\{e\}$ 恰好提示不能删除这个假设。

Frattini 论证要求 $N\lhd G$ 且 $P$ 为 $N$ 的 Sylow 子群，才能断言
$G=N N_G(P)$。
取 $N=R$、$P=R$ 时，确有
\[
G=R N_G(R)=RG.
\]
这是一个成立但退化的实例。
若误取 $N=R$、$P=H_0$，则 $P$ 根本不在 $N$ 内，不能套用该定理。

另外，\emph{Frattini 子群}的定义是所有极大子群的交，不能与上述论证混淆。
本例四个极大子群已经全部列出，故
\[
\Phi(G)=R\cap H_0\cap H_1\cap H_2=\{e\}.
\]
因此这里并无可隐藏在每个极大子群中的非单位生成元。
\end{workedexample}

\begin{workedexample}[同一个顶点作用中的轨道、稳定子和 Burnside 计数]
\label{ex:d3-comprehensive-actions}
在顶点集合 $X=\{1,2,3\}$ 上仍令 $r=(123)$、$s=(23)$。
求轨道、稳定子和作用的核，并计算 $q$ 种颜色的三角形在旋转、反射下有多少种不同着色。

\textbf{解。} 从顶点 $1$ 出发，旋转可以到达任一顶点。
轨道 $G\cdot1=X$，稳定子为 $G_1=\{e,s\}=H_0$，因此
\[
|G\cdot1|=3=6/2=|G|/|G_1|.
\]
此作用传递且忠实，但不自由：$s\ne e$ 仍固定顶点 $1$。
而限制到 $R$ 后，每个非单位旋转都不固定顶点，故 $R$ 对 $X$ 的作用自由、传递且正则。

左陪集作用 $G\curvearrowright G/H_0$ 与这个顶点作用相同：
\[
gH_0\longmapsto g\cdot1.
\]
其核是三个共轭稳定子的交 $H_0\cap H_1\cap H_2=\{e\}$。
对比 $G\curvearrowright G/R$：该作用传递，但核为 $R$，因此不忠实也不自由。
两个作用的差别可从稳定子直接读取，不需更换定义。

给每个顶点赋予 $q$ 种\emph{有名称}的颜色，颜色赋值共有 $q^3$ 个，
群只移动顶点，不置换颜色名称。
一个群元固定的赋值必须在其每个顶点循环上颜色相同，所以
\[
|\operatorname{Fix}(e)|=q^3,\qquad
|\operatorname{Fix}(r)|=|\operatorname{Fix}(r^2)|=q,\qquad
|\operatorname{Fix}(r^a s)|=q^2.
\]
Burnside 引理由此给出不计三角形转动或反射的着色数
\[
N(q)=\frac{q^3+2q+3q^2}{6}=\binom{q+2}{3}.
\]
例如 $q=2$ 时有 $4$ 类：按第一种颜色出现 $0,1,2,3$ 次区分。
分母是 $6$，但并不是每个着色轨道都有六个元素；单色赋值的稳定子是整个 $G$。
\end{workedexample}

\begin{workedexample}[Cayley 的六点作用与三点忠实作用的差别]
\label{ex:d3-comprehensive-cayley}
将集合 $G$ 的元素按 $(e,r,r^2,s,rs,r^2s)$ 编号为 $1,\ldots,6$。
写出左乘作用的生成元置换，并与三点顶点作用比较忠实性和自由性。

\textbf{解。} 按乘法表依次左乘每个群元，
左乘作用给出
\[
L_r=(123)(456),\qquad L_s=(14)(26)(35).
\]
例如 $L_s(r)=sr=r^2s$，对应 $2\mapsto6$，与乘法表一致。
若 $L_g$ 固定单位元，则 $g=L_g(e)=e$，所以这个作用自由且忠实；
任意 $a,b\in G$ 之间由 $ba^{-1}$ 左乘连接，所以它也传递，因而正则。
这就是 Cayley 嵌入 $G\hookrightarrow S_6$ 的具体实现。
三点作用给出了更小的忠实嵌入 $G\hookrightarrow S_3$，却不正则。
“忠实”只排除在\emph{所有点上}都恒等的非单位群元，不排除单个点有非平凡稳定子。
\end{workedexample}

\begin{workedexample}[合成列、可解性与中心列的停滞]
\label{ex:d3-comprehensive-series}
求合成列、导出列和上下中心列，判断 $G$ 是否可解、是否幂零。

\textbf{解。} 已知 $R$ 正规，链 $G\triangleright R\triangleright\{e\}$ 的因子为
\[
G/R\cong C_2,\qquad R/\{e\}\cong C_3.
\]
两个因子都是单群，故它是合成列。
链 $G>H_0>\{e\}$ 虽然每次指标都是素数，却不是次正规列：$H_0\not\lhd G$。
不能用它另造一套合成因子。

采用交换子 $[a,b]=aba^{-1}b^{-1}$，计算
\[
[s,r]=r^{-2}=r.
\]
因为 $G/R$ Abel，全部交换子都在 $R$ 内；又有交换子等于 $r$，所以
\[
G'=R,\qquad G''=R'=\{e\}.
\]
导出列在第二步到单位群，$G$ 因而可解，导出长度为 $2$。
这也解释了为什么所有一维表示都会杀掉 $R$。

但下中心列为
\[
\gamma_1(G)=G,\qquad\gamma_2(G)=[G,G]=R,
\qquad\gamma_n(G)=R\quad(n\ge2),
\]
因为 $[s,r]=r$ 已给出 $[G,R]=R$。
上中心列则从 $Z_0=\{e\}$ 出发，由 $Z(G)=\{e\}$ 得
\[
Z_1=Z_2=\cdots=\{e\}.
\]
两种中心列均不能到达相应终点，因此 $G$ 非幂零。
Sylow $2$-子群不正规也是\Cref{thm:nilpotent-equivalent}的有限群判据中的同一障碍。
\end{workedexample}

\begin{workedexample}[半直积、非直积与自由群展示的含义]
\label{ex:d3-comprehensive-semidirect}
用 $R,H_0$ 把 $G$ 写成半直积，并说明这个分解与原群展示的关系。

\textbf{解。} 前面已算出 $R\lhd G$、$H_0\cap R=\{e\}$、$RH_0=G$。
$H_0$ 在 $R$ 上的共轭作用为 $s:r\mapsto r^{-1}$，故
\[
G\cong C_3\rtimes C_2,
\qquad (a,b)(c,d)=(a+(-1)^b c,b+d).
\]
这正好是最初得到的乘法公式；半直积在这里不是另一种记号，而是重建了全部运算。
此作用非平凡，所以该分解不是 $C_3\times C_2$：后者 Abel 且有六阶元，而 $G$ 两者均不具备。

若 $F(x,y)$ 是两个字母的自由群，展示的严格含义是
\[
G\cong F(x,y)\big/\langle\!\langle
x^3,\ y^2,\ yxyx
\rangle\!\rangle,
\qquad x\mapsto r,\quad y\mapsto s.
\]
双尖括号表示所列关系生成的\emph{正规闭包}，不是普通子群生成。
商群中词可化为六个正规形，顶点置换又保证这六个像互不相同，
所以所列关系既足够，也未引入额外塌缩。
这不表示 $G$ 是自由群；例如非单位元 $r$ 满足 $r^3=e$，而自由群没有这样的有限阶元。
若改用两个反射 $p=s,q=rs$，则
\[
p^2=q^2=e,\qquad pq=r^{-1},\qquad
G=\langle p,q\mid p^2=q^2=(pq)^3=e\rangle.
\]
这与原展示等价：由 $r=qp,s=p$ 可相互恢复生成元。
因此旋转加反射的展示与两个反射的 Coxeter 展示描述的是同一运算。
\end{workedexample}

\begin{workedexample}[从生成元矩阵构造全部不可约复表示]
\label{ex:d3-comprehensive-irreps}
求全部一维复表示，构造一个二维复表示并证明其不可约，计算各表示的核。

\textbf{解。} 以下表示空间均在 $\mathbb C$ 上，二维实矩阵也按复线性延拓。
一维表示中，标量交换，所以 $srs=r^{-1}$ 强迫 $\rho(r)=\rho(r)^{-1}$；
再与 $\rho(r)^3=1$ 联用，得到 $\rho(r)=1$。
因 $\rho(s)^2=1$，全部一维表示恰为
\[
A_1(r)=A_1(s)=1,\qquad A_2(r)=1,\quad A_2(s)=-1.
\]
它们分别是平凡表示与刚才的同态 $\varepsilon$。

在平面 $E$ 的正交基 $x,y$ 上取
\[
D_E(r)=\begin{pmatrix}-\frac12&-\frac{\sqrt3}{2}\\
\frac{\sqrt3}{2}&-\frac12\end{pmatrix},\qquad
D_E(s)=\begin{pmatrix}1&0\\0&-1\end{pmatrix}.
\]
矩阵满足 $D_E(r)^3=D_E(s)^2=I$、$D_E(s)D_E(r)D_E(s)=D_E(r)^{-1}$，
所以展示保证它们确实定义表示。
令 $\omega=e^{2\pi i/3}$。$D_E(r)$ 的本征向量
\[
x-iy,\quad x+iy
\]
分别对应 $\omega,\omega^{-1}$，而 $D_E(s)$ 互换两条本征直线。
若有一维不变子空间，它必须是其中一条，却不能同时对 $s$ 不变；故 $E$ 在复数域上仍不可约。
$A_1$ 的核为 $G$，$A_2$ 的核为 $R$，$E$ 的核为 $\{e\}$。
不可约性与忠实性因此是不同问题：这里两个一维不可约表示均不忠实。
\end{workedexample}

\begin{workedexample}[酉化与 Schur 引理都化成一个二阶矩阵计算]
\label{ex:d3-comprehensive-schur}
求 $E$ 的自交织算子，以及 $E$ 与两个一维表示之间的交织算子；将 $E$ 换到非正交基后，再用群平均求不变内积。

\textbf{解。} $D_E(r),D_E(s)$ 本已正交，所以 $E$ 是酉表示。
若 $T$ 同时与二者交换，先与 $D_E(s)$ 交换得 $T=\operatorname{diag}(a,b)$；
再比较 $TD_E(r)$ 与 $D_E(r)T$ 的非对角元，得到 $a=b$。
因此
\[
\operatorname{End}_G(E)=\mathbb C I,
\]
这正是\Cref{thm:schur-lemma}对 $E$ 的结论。
从 $A_1$ 或 $A_2$ 到 $E$ 的交织算子若非零，其像将是一维不变子空间；
从 $E$ 到一维表示的交织算子若非零，秩为一，因而其核也将是一维不变子空间。
两者均与 $E$ 的不可约性矛盾；$A_1$ 与 $A_2$ 之间的交织条件在 $s$ 上又给出 $T=-T$，所以不同不可约表示之间的交织算子全为零。

也可故意换用不正交基：取 $T_0=\operatorname{diag}(2,1)$，令
$D'(g)=T_0^{-1}D_E(g)T_0$。
它仍是同一个表示，却一般不满足 $D'(g)^\dagger D'(g)=I$。
群平均给出的正定矩阵为
\[
Q=\frac16\sum_g D'(g)^\dagger D'(g)
=\frac58\,T_0^\dagger T_0
=\begin{pmatrix}5/2&0\\0&5/8\end{pmatrix}.
\]
这里令 $B=(T_0^{-1})^\dagger T_0^{-1}=\operatorname{diag}(1/4,1)$。
平均矩阵 $\frac16\sum_gD_E(g)^\dagger B D_E(g)$ 与群作用交换，故由刚求出的交换子代数可写成 $cI$。
共轭不改变迹，取迹得到 $2c=\operatorname{Tr}B=5/4$，这就确定了上式的 $5/8$。
因而 $D'(g)^\dagger QD'(g)=Q$，以 $Q^{1/2}$ 换基便恢复酉表示。
这也说明“某组矩阵不酉”不等于“该有限群表示不可酉化”。
\end{workedexample}

\begin{workedexample}[特征标的行列正交与完整性]
\label{ex:d3-comprehensive-character}
从已构造的表示求特征标表，检查行列正交和穷尽性，再计算实型指标与共轭类和的作用。

\textbf{解。} 由前面实际构造的矩阵取迹，按 $e$、旋转类、反射类的次序得到
\[
\begin{array}{c|rrr}
 &e&r,r^2&s,rs,r^2s\\\hline
\text{类大小}&1&2&3\\
\chi_{A_1}&1&1&1\\
\chi_{A_2}&1&1&-1\\
\chi_E&2&-1&0
\end{array}
\]
正交内积必须使用类大小：
\begin{align*}
\langle\chi_{A_1},\chi_{A_2}\rangle&=(1+2-3)/6=0,\\
\langle\chi_{A_1},\chi_E\rangle
=\langle\chi_{A_2},\chi_E\rangle&=(2-2+0)/6=0,\\
\|\chi_{A_1}\|^2=\|\chi_{A_2}\|^2&=(1+2+3)/6=1,\\
\|\chi_E\|^2&=(4+2+0)/6=1.
\end{align*}
全部共轭类只有三个，由\Cref{thm:character-basis}，这三个互不等价不可约表示已经穷尽，
并且 $1^2+1^2+2^2=6$。
不能仅因矩阵维数小，就宣称没有其他不可约表示；穷尽性用到的是这个理论判据。

列正交也能逐项核对：三列的无权平方和为
\[
1+1+4=6=|C_G(e)|,\quad
1+1+1=3=|C_G(r)|,\quad
1+1+0=2=|C_G(s)|.
\]
不同列的内积分别为 $1+1-2=0$、$1-1+0=0$、$1-1+0=0$。
这些中心化子大小正是前面从乘法求出的数，因此表示计算回接到了群结构。
$E$ 的实正交矩阵还给出 $E\cong E^*$：不变双线性形式
$B(z,w)=z_xw_x+z_yw_y$ 将向量识别为对偶向量。
这里是双线性形式，不能把复共轭误放进识别 $E\to E^*$ 的复线性映射。
所有特征标值均为实数也与 $g$ 和 $g^{-1}$ 在同一共轭类相符。
再实际计算 Frobenius--Schur 指标：两个旋转平方后仍在旋转类，三个反射平方后都是单位元，故
\[
\nu(E)=\frac16\sum_g\chi_E(g^2)=\frac{2+2(-1)+3\cdot2}{6}=1,
\qquad \nu(A_1)=\nu(A_2)=1.
\]
这验证三个不可约表示都是实型；它将与后面 $Q_8$ 的二维四元数型指标 $-1$ 形成区别。

最后把 Burnside 证明使用的类和也算出。令
$K_r=r+r^2$、$K_s=s+rs+r^2s$，它们在三个不可约表示上的标量分别为
\[
\begin{array}{c|rrr}
 &A_1&A_2&E\\\hline
K_r&2&2&-1\\
K_s&3&-3&0
\end{array}
\]
例如 $D_E(r)+D_E(r)^2=-I_2$，三面反射矩阵之和为零。
一般公式 $|C|\chi(g)/d$ 在这里给出整数，因此类和的代数整性已有直接计算的实例。
\end{workedexample}

\begin{workedexample}[群代数、中央幂等元与六维正则表示]
\label{ex:d3-comprehensive-regular}
分解六维正则表示，求中央幂等元，并为每个不可约副本构造一组实际基向量。

\textbf{解。} 群代数 $\mathbb C[G]$ 以六个群元为正交基，左乘 $L_g$ 就是 Cayley 例题的置换矩阵。
故 $\chi_{\mathrm{reg}}=(6,0,0)$；重数公式依次给出 $1,1,2$，于是
\[
\mathbb C[G]\cong A_1\oplus A_2\oplus E\oplus E.
\]
其中两个 $E$ 是不同副本，不能误写成一个四维不可约表示。

把特征标投影公式写成群代数中央元，得到
\begin{align*}
z_{A_1}&=\frac16(e+r+r^2+s+rs+r^2s),\\
z_{A_2}&=\frac16(e+r+r^2-s-rs-r^2s),\\
z_E&=\frac13(2e-r-r^2).
\end{align*}
例如令 $a=e+r+r^2$、$b=s+rs+r^2s$，乘法表直接给出
\[
a^2=3a,\quad ab=ba=3b,\quad b^2=3a.
\]
所以 $z_{A_1}^2=z_{A_1}$、$z_{A_2}^2=z_{A_2}$、$z_{A_1}z_{A_2}=0$；
又 $z_E=e-z_{A_1}-z_{A_2}$，故三个幂等元两两正交且和为单位。
它们的系数在共轭类上相同，所以都中央。

为把四维 $E$ 等型空间再拆成两个副本，定义六个向量
\[
f_{\alpha,ij}=\sqrt{\frac{d_\alpha}{6}}
\sum_{g\in G}\overline{D_\alpha(g)_{ij}}\,g,
\quad \alpha\in\{A_1,A_2,E\},\quad1\le i,j\le d_\alpha.
\]
大正交关系给出它们标准正交。左乘 $h$ 后令 $g=h^{-1}k$，再用
$D_\alpha(h^{-1}k)=D_\alpha(h)^\dagger D_\alpha(k)$ 展开矩阵元，得到
\[
L_h f_{\alpha,ij}=\sum_{k=1}^{d_\alpha}
D_\alpha(h)_{ki}f_{\alpha,kj}.
\]
因此固定 $j$ 的子空间是一个 $\alpha$ 副本；$E$ 的 $j=1,2$ 给出两个二维副本。
这六个向量构成全空间的基，从而给出了正则分解的实际换基，
而 $L_{z_E}$ 本身只能投到两个副本的\emph{总和}。
\end{workedexample}

\begin{workedexample}[六点函数的 Fourier 展开与 Plancherel]
\label{ex:d3-comprehensive-fourier}
用不可约表示矩阵元展开任意六点函数 $f:G\to\mathbb C$，求反演和范数公式，并以单位元指示函数检验归一化。

\textbf{解。} 沿用特征标内积的归一化，定义
\[
\widehat f_{\alpha,ij}=\frac16\sum_{g\in G}
\overline{D_\alpha(g)_{ij}}f(g).
\]
上一例题的大正交关系和六个基向量的完备性直接给出
\[
f(g)=\sum_\alpha d_\alpha\sum_{i,j}
\widehat f_{\alpha,ij}D_\alpha(g)_{ij},\qquad
\frac16\sum_g|f(g)|^2
=\sum_\alpha d_\alpha\sum_{i,j}|\widehat f_{\alpha,ij}|^2.
\]
本例只用有限求和，不需要调用紧群的一般 Peter--Weyl 定理。
取 $f=\delta_e$，则 $\widehat f_{\alpha,ij}=\delta_{ij}/6$，因此
\[
\delta_e(g)=\frac16\bigl[\chi_{A_1}(g)+\chi_{A_2}(g)+2\chi_E(g)\bigr].
\]
右边在三类上分别为 $(1+1+4)/6=1$、$(1+1-2)/6=0$、$(1-1+0)/6=0$，
正好重建单位元指示函数。
其 Plancherel 等式为 $1/6=(1^2+1^2+2^2)/36$。
若 $f$ 是类函数，Schur 引理使各 Fourier 矩阵为标量矩阵；
因此一般的六点函数需要六个矩阵元系数，而类函数只需要三个特征标系数。
\end{workedexample}

\begin{workedexample}[三维顶点置换表示的实际投影与基底]
\label{ex:d3-comprehensive-permutation}
在 $V=\mathbb C^3$ 中取顶点基 $e_1,e_2,e_3$，令 $P(g)e_i=e_{g\cdot i}$。
求不可约分解、投影矩阵和对称适配基，并把任意顶点振幅分到这些模式中。

\textbf{解。} 特征标等于不动顶点数。
仍使用同一个 $r=(123)$、$s=(23)$，其特征标为不动顶点数：
\[
\chi_V=(3,0,1),\qquad
(m_{A_1},m_{A_2},m_E)=(1,0,1).
\]
将上一例题的中央元送入此表示，得到
\[
\Pi_{A_1}=\frac13
\begin{pmatrix}1&1&1\\1&1&1\\1&1&1\end{pmatrix},\qquad
\Pi_{A_2}=0,\qquad \Pi_E=I-\Pi_{A_1}.
\]
所以 $A_1$ 是所有顶点等幅的方向，$E$ 则是分量和为零的平面。
正交基可明确取为
\[
b_0=\frac{e_1+e_2+e_3}{\sqrt3},\quad
b_x=\frac{2e_1-e_2-e_3}{\sqrt6},\quad
b_y=\frac{e_2-e_3}{\sqrt2}.
\]
直接代入顶点置换得
\[
P(r)b_x=-\tfrac12 b_x+\tfrac{\sqrt3}{2}b_y,\qquad
P(r)b_y=-\tfrac{\sqrt3}{2}b_x-\tfrac12 b_y,
\]
而 $P(s)b_x=b_x$、$P(s)b_y=-b_y$。
故限制到 $\operatorname{span}\{b_x,b_y\}$ 的矩阵恰好是前面构造的 $D_E$，
不是仅有相同维数的一块。

若初始顶点振幅为 $u=(u_1,u_2,u_3)^T$，其对称与非对称分量分别为
\[
\Pi_{A_1}u=\bar u(1,1,1)^T,\qquad
\Pi_Eu=(u_1-\bar u,u_2-\bar u,u_3-\bar u)^T,
\quad \bar u=(u_1+u_2+u_3)/3.
\]
这一步使抽象的重数、等型空间和投影成为可直接计算的模式幅度。
例如先选一个未必正交的投影 $Q$ 到 $\mathbb C b_0$，
$Qw=\ell(w)b_0$、$\ell(b_0)=1$。
群平均后
\[
\frac16\sum_gP(g)QP(g)^{-1}w
=\left[\frac16\sum_g\ell(P(g)^{-1}w)\right]b_0
=\Pi_{A_1}w.
\]
最后的平均线性泛函对全部顶点置换不变，故必是各坐标等权之和，
并由 $b_0$ 上取值 $1$ 确定归一化。
于是 $\ker\Pi_{A_1}=E$ 是不变补空间，具体实现了 Maschke 的平均投影构造。
\end{workedexample}

\begin{workedexample}[限制、商群提升与两个子群的诱导矩阵]
\label{ex:d3-comprehensive-induction}
求三个不可约表示在 $R,H_0$ 上的限制，并从这两个子群的一维表示显式构造诱导矩阵与分解。

\textbf{解。} 先把已有矩阵限制到两个子群。
记 $R$ 的三个一维表示为 $\sigma_m(r)=\omega^m$，$m=0,1,2$，
记 $H_0$ 的两个一维表示为 $\tau_\pm(s)=\pm1$。
旋转本征基与反射本征基分别给出
\[
E\downarrow_R=\sigma_1\oplus\sigma_2,\qquad
E\downarrow_{H_0}=\tau_+\oplus\tau_-.
\]
$A_1,A_2$ 限制到 $R$ 均为 $\sigma_0$，限制到 $H_0$ 分别为 $\tau_+,\tau_-$。
因为 $R$ 位于 $A_1,A_2$ 的核内，它们恰是 $G/R=C_2$ 的两个表示通过商映射提升所得；
$E$ 在 $R$ 上非平凡，故不能下降为 $G/R$ 的表示。

诱导 $\sigma_m$ 时选左陪集代表 $e,s$，在基
$(e\otimes1,s\otimes1)$ 中逐项使用 $rs=sr^{-1}$，得到
\[
D_{\operatorname{Ind}_R^G\sigma_m}(r)
=\begin{pmatrix}\omega^m&0\\0&\omega^{-m}\end{pmatrix},\qquad
D_{\operatorname{Ind}_R^G\sigma_m}(s)
=\begin{pmatrix}0&1\\1&0\end{pmatrix}.
\]
当 $m=1$ 时，映射 $e\otimes1\mapsto(x-iy)/\sqrt2$、
$s\otimes1\mapsto(x+iy)/\sqrt2$ 是到 $E$ 的酉交织映射。
当 $m=2$ 时交换两基矢即得同一表示。
当 $m=0$ 时两矩阵分别为 $I$ 与交换矩阵，其对称、反对称方向给出
\[
\operatorname{Ind}_R^G\sigma_0=A_1\oplus A_2,\qquad
\operatorname{Ind}_R^G\sigma_1\cong\operatorname{Ind}_R^G\sigma_2\cong E.
\]

诱导 $\tau_\pm$ 时选代表 $e,r,r^2$，令 $c_a=r^a\otimes1$。
生成元的作用为
\[
r:c_a\mapsto c_{a+1},\qquad
s:c_a\mapsto\tau_\pm(s)c_{-a}.
\]
所以特征标分别为 $(3,0,1)$ 与 $(3,0,-1)$，并由重数公式得
\[
\operatorname{Ind}_{H_0}^G\tau_+=A_1\oplus E,\qquad
\operatorname{Ind}_{H_0}^G\tau_-=A_2\oplus E.
\]
前者就是顶点置换表示；后者在其基础上给反射乘一个负号。
在二维补空间中，取
\[
J=\begin{pmatrix}0&-1\\1&0\end{pmatrix},\qquad
JD_E(g)J^{-1}=\varepsilon(g)D_E(g),
\]
因为 $J$ 与旋转交换、与反射反交换。
这个交织矩阵消去反射负号，故二维成分仍等价于 $E$。
\end{workedexample}

\begin{workedexample}[Frobenius、Mackey 与 Clifford 的同一组数据]
\label{ex:d3-comprehensive-clifford}
利用上一例题的诱导矩阵，逐项检验 Frobenius 互反律和 Mackey 限制公式，并求正规子群 $R$ 的字符轨道与惯性群。

\textbf{解。} Frobenius 互反律把限制与诱导中的重数对应起来。
例如 $E\downarrow_R$ 含 $\sigma_1$ 一次，对应
$\operatorname{Ind}_R^G\sigma_1=E$ 中 $E$ 一次；
$A_1\downarrow_R=A_2\downarrow_R=\sigma_0$，对应
$\operatorname{Ind}_R^G\sigma_0=A_1\oplus A_2$。
对 $H_0$，$E$ 的两个反射本征值各出现一次，
对应两种 $\operatorname{Ind}_{H_0}^G\tau_\pm$ 都含一个 $E$。
两边的重数都为一，与限制矩阵的本征直线数一致。

正规子群 $R$ 的双陪集为 $R\backslash G/R=\{R,RsR\}$，代表元取 $e,s$。
两项的交子群都为 $R$，且共轭表示为
\[
{}^s\sigma_m(r)=\sigma_m(s^{-1}rs)=\omega^{-m}.
\]
Mackey 公式于是变成
\[
(\operatorname{Ind}_R^G\sigma_m)\downarrow_R
\cong\sigma_m\oplus\sigma_{-m},
\]
恰好恢复前一例题的两个对角矩阵元。
对 $H_0$，双陪集代表可取 $e,r$，因为
\[
H_0rH_0=\{r,r^2,rs,r^2s\},\qquad
H_0\cap rH_0r^{-1}=\{e\}.
\]
故第二项是从单位子群诱导的 $H_0$ 正则表示，得到
\[
(\operatorname{Ind}_{H_0}^G\tau_\pm)\downarrow_{H_0}
\cong\tau_\pm\oplus\tau_+\oplus\tau_-.
\]
其 $s$ 上的迹是 $\pm1+1-1=\pm1$，与三维诱导矩阵完全一致。

Clifford 定理研究\emph{正规}子群限制。
$G$ 对 $\widehat R=\{\sigma_0,\sigma_1,\sigma_2\}$ 的轨道是
\[
\{\sigma_0\},\qquad\{\sigma_1,\sigma_2\}.
\]
$\sigma_0$ 的惯性群为 $G$，$\sigma_1,\sigma_2$ 的惯性群为 $R$。
$A_1,A_2$ 对应第一轨道且限制重数为 $1$，$E$ 对应第二轨道且两项各重数 $1$，
维数公式分别为 $1=1\cdot1\cdot1$ 与 $2=1\cdot2\cdot1$。
在半直积的小群语言中，第一轨道的反射稳定子为 $C_2$，其两个表示产生 $A_1,A_2$；
第二轨道的反射稳定子是单位群，其唯一表示诱导产生 $E$。
本例无需射影表示障碍，因为两个稳定子上的延拓均已由明确矩阵实现。
但 $H_0$ 非正规，不能把 $E\downarrow_{H_0}$ 的两项强称为 Clifford 正规子群轨道。

最后，$|G|=2\cdot3$ 满足 Burnside $p^aq^b$ 定理的阶数条件，
其可解结论与前面独立算出的 $G'=R,G''=\{e\}$ 相符。
旋转类大小 $2$、反射类大小 $3$ 也都是素数幂；
本例确实存在非平凡真正规子群 $R$，所以并非非交换单群。
这些都是一般结论的具体实例，六阶计算并不能代替这些一般定理的证明。
\end{workedexample}

\begin{workedexample}[张量平方的 Clebsch--Gordan 系数与交换对称性]
\label{ex:d3-comprehensive-cg}
分解 $E\otimes E$，求耦合基、投影和对称平方；据此写出全部不可约类型的张量乘法。

\textbf{解。} 在 $E\otimes E$ 的有序基 $(xx,xy,yx,yy)$ 中，$xy$ 表示 $x\otimes y$。
张量积特征标为 $\chi_{E\otimes E}=\chi_E^2=(4,1,0)$，重数计算给出
\[
m_{A_1}=m_{A_2}=(4+2)/6=1,\qquad m_E=(8-2)/6=1,
\]
所以 $E\otimes E\cong A_1\oplus A_2\oplus E$。
其余张量积也能在同一模型内补全：$A_1\otimes\alpha\cong\alpha$、
$A_2\otimes A_2\cong A_1$，而 $A_2\otimes E\cong E$。
最后一个等价正由诱导例题中已算出的交织矩阵 $J$ 实现。
因此表示环以 $1=[A_1]$、$a=[A_2]$、$v=[E]$ 为整数加法基，乘法由
\[
a^2=1,\qquad av=v,\qquad v^2=1+a+v
\]
完全确定。这里的整数线性组合允许虚表示，例如 $[V]-[A_1]=[E]$；
实际表示的不可约重数则必须非负，不能把任意环元素直接称为态空间。
但特征标还没给出耦合基。明确取
\[
t=\frac{xx+yy}{\sqrt2},\quad
a=\frac{xy-yx}{\sqrt2},\quad
u=\frac{xx-yy}{\sqrt2},\quad
v=-\frac{xy+yx}{\sqrt2}.
\]
四个向量标准正交，其列组成的换基矩阵为
\[
C=\frac1{\sqrt2}
\begin{pmatrix}
1&0&1&0\\0&1&0&-1\\0&-1&0&-1\\1&0&-1&0
\end{pmatrix},\qquad C^\dagger C=I_4.
\]
二阶正交变换保持 $xx+yy$，故 $t$ 按 $A_1$ 变换；
反对称张量 $a$ 按行列式变换，故 $r:a\mapsto a$、$s:a\mapsto-a$，是 $A_2$。
若平面旋转角为 $\theta$，$(u,v)$ 上的矩阵为旋转 $-2\theta$；
在本例 $\theta=2\pi/3$ 时，$-2\theta\equiv\theta\pmod{2\pi}$。
反射则仍给出 $u\mapsto u,v\mapsto-v$，因此
\[
C^\dagger[D_E(g)\otimes D_E(g)]C
=\operatorname{diag}\bigl(A_1(g),A_2(g),D_E(g)\bigr)
\quad(g\in G).
\]
检验只需对生成元 $r,s$ 做，因为其余群元都是它们的词。

在这个空间中中央投影于是具体等于
\[
\Pi_{A_1}=tt^\dagger,\qquad
\Pi_{A_2}=aa^\dagger,\qquad
\Pi_E=uu^\dagger+vv^\dagger.
\]
对一般张量 $T=T_{xx}xx+T_{xy}xy+T_{yx}yx+T_{yy}yy$，耦合系数就是
\[
\frac1{\sqrt2}(T_{xx}+T_{yy},\ T_{xy}-T_{yx},\ T_{xx}-T_{yy},\ -T_{xy}-T_{yx})^T.
\]
交换两因子的算子固定 $t,u,v$ 而使 $a$ 变号，所以
\[
\operatorname{Sym}^2 E=A_1\oplus E,\qquad\Lambda^2E=A_2.
\]
尤其相同向量的平方 $w\otimes w$ 没有 $A_2$ 分量；这里“未出现”的原因是交换对称性，
不是因为 $A_2$ 在群中不存在。
以交织空间的语言说，$\dim\operatorname{Hom}_G(E\otimes E,\alpha)=1$
对 $\alpha=A_1,A_2,E$ 均成立，三个通道分别由上述迹、反对称和无迹对称分量给出。
因此仅凭 $D_3$ 对称性不禁止 $E\otimes E\to A_2$；
增加相同变量的对称交换要求后，该通道才消失。
\end{workedexample}

\begin{workedexample}[把群结构与耦合语言用于一个完整的动力学模型]
\label{ex:d3-comprehensive-model}
考虑三角形三个顶点各一个标量位移 $q_i$，每个顶点质量为 $m>0$，采用刚才的置换表示。
求 $D_3$ 对称的实对称刚度矩阵、正则模式和频率，再代入等强顶点差弹簧模型。

\textbf{解。} 顶点质量相同，动能在原顶点基中为 $\tfrac m2\sum_i\dot q_i^2$。
若实对称刚度矩阵 $K$ 与所有 $P(g)$ 对易，顶点置换迫使三个对角元相等、三个非对角元相等，故
\[
K=\begin{pmatrix}a&b&b\\b&a&b\\b&b&a\end{pmatrix}
=(a+2b)\Pi_{A_1}+(a-b)\Pi_E.
\]
在 $(b_0,b_x,b_y)$ 基中，势能与动能分别成为
\[
U=\tfrac12(a+2b)Q_0^2+\tfrac12(a-b)(Q_x^2+Q_y^2),\qquad
T=\tfrac m2(\dot Q_0^2+\dot Q_x^2+\dot Q_y^2).
\]
若 $a+2b>0$、$a-b>0$，频率为
\[
\omega_{A_1}^2=(a+2b)/m,\qquad
\omega_E^2=(a-b)/m,
\]
后者必二重简并。这是 $E$ 上 Schur 引理的动力学体现。
若取等强顶点差弹簧
\[
U=\frac k2\bigl[(q_1-q_2)^2+(q_2-q_3)^2+(q_3-q_1)^2\bigr],
\]
则 $a=2k,b=-k$，$A_1$ 是零频整体位移，$E$ 的两模式均有 $\omega^2=3k/m$。
这里标量位移模型有三个自由度，并非三个空间矢量位移的九维分子模型，不能混用模式数。

最后，二次模式组合由上一例题直接分类：
\[
Q_x^2+Q_y^2\sim A_1,\qquad
(Q_x^2-Q_y^2,-2Q_xQ_y)\sim E.
\]
这给出不变量与非平凡耦合的具体公式。
至此，从生成关系到陪集、群作用，再到不可约表示、投影、简并和张量耦合，
所有计算仍围绕同一个 $D_3$ 模型，并可回到最初的六元素乘法表逐步核查。
\end{workedexample}
